\documentclass[11pt,a4paper]{article}
\usepackage[T1]{fontenc}
\usepackage[utf8]{inputenc}
\usepackage[main=italian,provide=*]{babel}
\usepackage{amsmath,amssymb}
\usepackage{geometry}
\geometry{margin=2.5cm}
\usepackage{booktabs}
\usepackage{seqsplit}
\usepackage{hyperref}
\hypersetup{colorlinks=true,linkcolor=blue,citecolor=blue,urlcolor=blue}

\newcommand{\ket}[1]{\lvert #1 \rangle}

\title{Correlazioni dinamiche — trimero ad anello}
\author{Note di lavoro — Tesi triennale, Samuele Schiavi\\ Relatore: Prof.\ Alessandro Chiesa}
\date{}

\begin{document}
\maketitle

\section*{Introduzione}

Questo documento raccoglie i file della fase di \emph{correlazioni dinamiche}
(circuito con l'ancilla, Hadamard test) per il trimero ad anello — mirror diretto,
per $N=3$, di quanto già completato per il dimero
(\texttt{\seqsplit{circuito\_correlatori\_spiegato.tex}},
\texttt{\seqsplit{circuito\_correlazioni\_tutte.ipynb}}). A differenza del dimero,
qui l'analisi di simmetria è stata condotta e verificata numericamente
\textbf{prima} di scrivere una riga di circuito — lezione esplicita imparata sui due
bug del circuito dimero (coniugato complesso mancante nella formula spettrale,
mappa sito$\to$qubit invertita), entrambi mascherati dalla simmetria $U$ sui primi
correlatori testati e visibili solo dopo lo scan sistematico. Per il trimero
l'ordine è stato invertito apposta.

Per ciascun file: cosa fa, le funzioni interne, come si usa, le verifiche
effettuate nel corso della fase, e cosa se ne può concludere.

\section*{\texttt{\seqsplit{simmetrie\_correlatori\_trimero\_anello.tex}}}

Documento di teoria, primo passo della fase: analisi di simmetria dell'Hamiltoniana
completa (scambio + campo + DM Opzione B) applicata ai correlatori
$C_{ij}^{\alpha\beta}(t)=\langle\psi_0|\sigma_i^\alpha(t)\sigma_j^\beta(0)|\psi_0\rangle$,
prima di qualunque implementazione di circuito.

\textbf{Contenuto.} Lo scambio puro $P_{12}$ resta rotto dal DM (dispari sotto
scambio, come nel dimero); la generalizzazione della riparazione del dimero non è
banale — aggiungere l'identità sul terzo qubit \textbf{non} funziona
($\|[U_\text{naive},H]\|\sim31$ su parametri casuali, rottura netta, non un
residuo piccolo). La simmetria corretta richiede $R_z(\pi)$ su tutti e tre i qubit:
$$U_\text{anello} = \mathrm{SWAP}_{12}\cdot\big(R_z(\pi)\big)^{\otimes 3},$$
verificata $[U_\text{anello},H]=0$ a precisione macchina su 20 punti casuali
$(J,J',b,D)$. Differenza qualitativa dal dimero: $U_\text{anello}^2=-\mathbb I$
(ordine 4, non 2), ma resta unitaria con autovalori di modulo 1 — sufficiente per
l'argomento sui correlatori. Lemma:
$U_\text{anello}\,\sigma_i^\alpha\,U_\text{anello}^\dagger=\eta_\alpha\,
\sigma_{\pi(i)}^\alpha$, $\pi=(1\,2)$ (fissa il sito 3), $\eta_x=\eta_y=-1,\eta_z=+1$.
\textbf{Corollario nuovo, specifico di $N=3$} (non presente nel dimero, che non ha
un sito fisso sotto lo scambio): il sito 3 è punto fisso di $\pi$, quindi
$C_{33}^{\alpha\beta}(t)=\eta_\alpha\eta_\beta\,C_{33}^{\alpha\beta}(t)$ per ogni
$t$ — zero rigoroso per ogni $t$ quando $\eta_\alpha\eta_\beta=-1$:
$C_{33}^{xz}=C_{33}^{zx}=C_{33}^{yz}=C_{33}^{zy}\equiv0$. Simmetria di
time-reversal $K$ (stesso argomento del dimero, $H$ reale in base computazionale
per ogni $J,J',b,D$): zero rigoroso a $t=0$ per ogni coppia $i\neq j$ con
esattamente una componente $y$.

\textbf{Verifiche.} Scan classico completo delle 81 combinazioni al punto di
lavoro confermato ($J{=}1,J'{=}0.4,b{=}b_c{=}2.4,D{=}0.15$, Opzione B): esattamente
4 nulle per ogni $t$, nessuno zero aggiuntivo non spiegato da simmetria. Compilato
(\texttt{pdflatex}, due passate), zero errori/overfull, 8 pagine.

\textbf{Conclusione.} Punto di lavoro confermato a posteriori (non solo assunto);
nessuno zero strutturale nascosto — via libera per implementare il circuito sapendo
già cosa aspettarsi.

\section*{\texttt{\seqsplit{circuito\_correlazioni\_trimero\_anello.py}}}

Modulo di produzione principale: costruisce il circuito Hadamard test a 4 qubit (3
di registro + 1 ancilla) per misurare $C_{ij}^{\alpha\beta}(t)$, mirror diretto del
circuito già validato per il dimero, riusando
\texttt{\seqsplit{trotter\_trimero\_anello.py}} come blocco $U(t)$ — nessuna
ridefinizione indipendente dell'evoluzione.

\textbf{Funzioni interne.}
\begin{itemize}
\item \texttt{ANCILLA = 3} — indice del qubit ancilla (i 3 qubit di registro sono
0,1,2, stessa convenzione sito$\to$qubit di \texttt{\seqsplit{trimer\_ring\_exact.py}}).
\item \texttt{PAULI\_GATE} — dizionario componente$\to$gate ($X$,$Y$,$Z$), usato sia
per il controlled-$W$ sia per l'anti-controlled-$V$.
\item \texttt{ground\_state(J,\ Jp,\ b,\ D)} — stato fondamentale esatto (fase reale
fissata), da \texttt{\seqsplit{trimer\_ring\_exact.py}}.
\item \texttt{build\_correlator\_circuit(i,\ alpha,\ j,\ beta,\ t,\ N,\ J,\ Jp,\ b,\
D,\ part,\ psi0,\ measure)} — il circuito completo: \texttt{prepare\_state} sul
registro, $H$ sull'ancilla, controlled-$W=\sigma_j^\beta$ prima di $U(t)$, $U(t)$
Trotter \textbf{non} controllato, anti-controlled-$V=\sigma_i^\alpha$ dopo,
rotazione finale ($H$ per Re, $R_x(\pi/2)$ per Im) ed eventuale misura.
\item \texttt{ancilla\_z\_expectation(qc)} — lettura esatta di
$\langle\sigma_z^{anc}\rangle$ da statevector (nessun rumore di shot).
\item \texttt{correlator\_from\_circuit(...)} — esegue entrambe le parti (Re/Im) e
combina il risultato complesso.
\end{itemize}

\textbf{Come si usa.} Libreria, non notebook: si importano le funzioni necessarie.
Richiede \texttt{\seqsplit{trimer\_ring\_exact.py}} e
\texttt{\seqsplit{trotter\_trimero\_anello.py}} nella stessa cartella.

\textbf{Verifiche.} Nessun bug trovato in questa fase (a differenza del dimero) —
l'analisi di simmetria precedente aveva già eliminato i due tipi di errore più
insidiosi (segno mancante, siti scambiati) prima ancora di scrivere il circuito.

\textbf{Conclusione.} Modulo maturo, riusato senza modifiche da tutti gli script e
i notebook successivi della fase.

\section*{\texttt{\seqsplit{validate\_circuito\_correlazioni.py}}}

Modulo di validazione indipendente: riferimento classico esatto per
$C_{ij}^{\alpha\beta}(t)$, calcolato deliberatamente per esponenziale di matrice
diretto, non tramite la formula spettrale $\sum_k a_kb_k$ — scelta per non poter
reintrodurre per disattenzione la stessa classe di bug del "coniugato complesso
mancante" già trovata nel dimero (qui il confronto è fra due calcoli
concettualmente indipendenti, non fra una formula e la sua stessa riscrittura).

\textbf{Funzioni interne.}
\begin{itemize}
\item \texttt{I2, X, Y, Z, PAULI} — matrici di Pauli $2\times2$ e dizionario
componente$\to$matrice.
\item \texttt{kron3(A,\ B,\ C)} — prodotto di Kronecker a tre fattori.
\item \texttt{site\_op(site,\ alpha)} — operatore di Pauli a un sito, esteso a
$8\times8$ con identità sugli altri due qubit.
\item \texttt{classical\_exact(i,\ alpha,\ j,\ beta,\ t,\ J,\ Jp,\ b,\ D,\ psi0,\
H)} — $\langle\psi_0|e^{iHt}\sigma_i^\alpha e^{-iHt}\sigma_j^\beta|\psi_0\rangle$
via \texttt{scipy.linalg.expm}.
\end{itemize}

\textbf{Come si usa.} Libreria; \texttt{site\_op} e \texttt{classical\_exact} sono
importate ovunque nella fase (script e notebook) come riferimento di validazione.

\textbf{Verifiche.} Concordanza con \texttt{correlator\_from\_circuit} entro
l'errore di Trotter atteso, su tutti i casi testati nella fase (6 rappresentativi,
poi tutte le 81 combinazioni — vedi sotto).

\textbf{Conclusione.} Riferimento indipendente stabile, mai modificato dopo la
prima stesura: nessuna necessità emersa nel corso della fase.

\section*{\texttt{\seqsplit{circuito\_correlazioni\_trimero\_anello\_spiegato.tex}}}

Documento pedagogico principale: spiegazione passo-passo del circuito (stato di
partenza, perché il correlatore non è misurabile direttamente, costruzione stadio
per stadio, dimostrazione che l'ancilla legge l'overlap, i due punti delicati
$U(t)$ non controllato e anti-controllo su $V$, struttura del blocco Trotter,
misura, validazione finale), mirror per $N=3$ di
\texttt{\seqsplit{circuito\_correlatori\_spiegato.tex}} (dimero) — non una copia:
due sezioni riscritte da zero perché $N=3$ introduce differenze reali.

\textbf{Contenuto specifico di $N=3$, non presente nel mirror del dimero.}
(1) \emph{Il caso interessante: quando lo spalmamento non succede} — usa lo zero
strutturale del sito 3 (punto fisso di $U_\text{anello}$, senza equivalente nel
dimero) come esempio guida, invece di ripetere l'espansione di Heisenberg su tutti
i $4^3=64$ prodotti di Pauli. (2) \emph{Il blocco $U(t)$: Trotter a due livelli
annidati} — $H_{ex}$ e $H_{DM}$, somme di tre legami che condividono qubit a due a
due (la chiusura ad anello), richiedono ciascuno un Trotter interno oltre al
Trotter esterno già presente nel dimero.

\textbf{Aggiornamenti successivi (stessa fase, sotto-sezioni aggiunte in corso
d'opera).} "Validazione a shot finiti: quale errore domina?" — confronto Trotter
vs statistico, conclusione opposta al dimero (qui lo statistico domina di
$\sim16\times$, non il Trotter). "Scan sistematico delle 81 combinazioni, via
circuito effettivo" (Sez.~8.1) — i quattro zeri strutturali confermati via
circuito reale (non solo classicamente), tabella degli errori aggregati, la
scoperta di $C_{33}^{zz}\approx0.9995$ (autocorrelazione quasi congelata, fatto
dinamico non implicato da alcuna simmetria trovata).

\textbf{Figure prodotte} (da
\texttt{\seqsplit{generate\_figures\_circuito\_trimero\_anello.py}},
\texttt{\seqsplit{generate\_circuit\_fig\_trimero\_anello.py}},
\texttt{\seqsplit{generate\_figure\_shotnoise.py}},
\texttt{\seqsplit{generate\_figure\_scan81.py}} — dati dalla validazione
effettivamente svolta, non placeholder): diagramma del circuito, convergenza
Trotter, validazione continua $C_{11}^{yy}(t)$, zero strutturale del sito 3 su
intervallo continuo di $t$, convergenza vs shot, heatmap $9\times9$ delle 81
combinazioni.

\textbf{Verifiche.} Compilato più volte nel corso della fase (\texttt{pdflatex},
due passate ad ogni aggiornamento), zero errori; un solo underfull hbox cosmetico
residuo (badness 1735), invariato attraverso tutti gli aggiornamenti; due errori
di font expansion (\texttt{microtype}+\texttt{tcolorbox}) e due spezzature di
parola risolti alla prima stesura. 13 pagine alla versione più recente (cresciuto
da 10 attraverso tre aggiornamenti). \textbf{Nota:} da quando è stata introdotta la
regola "solo sorgente \texttt{.tex}, niente ricompilazione PDF" (per non consumare
token), le sotto-sezioni più recenti non sono state riverificate visivamente pagina
per pagina come le precedenti — il contenuto è comunque tratto direttamente dai
risultati numerici già verificati negli script corrispondenti.

\textbf{Conclusione.} Documento di riferimento per la difesa: ogni affermazione
della fase applicativa vi si può ricondurre con un riferimento di sezione preciso.

\section*{\texttt{\seqsplit{validate\_shot\_noise\_trimero\_anello.py}} + \texttt{\seqsplit{generate\_figure\_shotnoise.py}}}

Prima estensione a rumore statistico finito (mirror di
\texttt{\seqsplit{circuito\_correlazioni\_tutte.ipynb}} del dimero): fino a questo
punto il circuito era stato validato solo su statevector esatto.

\textbf{Contenuto.} \texttt{AerSimulator()} (shot-based); stima
$\langle\sigma_z^{anc}\rangle=(n_0-n_1)/N_\text{shots}$ da conteggi, per le due
esecuzioni separate (Re/Im) di ciascun correlatore. Nota tecnica: un circuito con
\texttt{prepare\_state} va transpilato (\texttt{transpile(qc,\ backend)}) prima di
\texttt{AerSimulator.run()}, altrimenti \texttt{AerError:\ unknown\ instruction:\
state\_preparation}. Tre correlatori rappresentativi testati
($C_{11}^{yy}$, $C_{33}^{xz}\approx0$, $C_{12}^{yy}$); convergenza vs
$N_\text{shots}$ ($256\to16384$, 40 ripetizioni indipendenti) verificata
$\sim1/\sqrt{N_\text{shots}}$.

\textbf{Risultato quantitativo centrale.} Errore di Trotter ($N=200$) $=7.08\times
10^{-4}$; soglia statistica nel caso peggiore $1/\sqrt{8192}=0.0110$; rapporto
statistico/Trotter $\sim15.6\times$ — \textbf{l'opposto} del regime dimostrativo
$R_1$ del dimero, dove il Trotter dominava di oltre $30\times$.

\textbf{Come si usa.} Script standalone, $\sim44$ s di esecuzione (cache di
transpilazione).

\textbf{Verifiche.} Nessuna anomalia; l'andamento $1/\sqrt{N_\text{shots}}$ atteso
è confermato quantitativamente, non solo qualitativamente.

\textbf{Conclusione.} A questo punto di lavoro, con $8192$ shots, non ha senso
spingere $N$ (Trotter) molto oltre 200 senza aumentare anche gli shots — il collo
di bottiglia è il rumore di campionamento, non l'approssimazione dell'evoluzione.

\section*{\texttt{\seqsplit{scan81\_trimero\_anello.py}} + \texttt{\seqsplit{generate\_figure\_scan81.py}}}

Estensione a \textbf{tutte e 81} le combinazioni $(i,j,\alpha,\beta)$ dello scan
finora fatto solo su 6 casi (statevector) e 3 casi (shot finiti) — misurando
ciascuna via circuito effettivo (non più solo classicamente), sia a statevector
esatto sia a shot finiti ($8192$ shots), allo stesso punto di lavoro confermato,
$t=1.3$, $N=200$. Lo scan a 81 combinazioni via \emph{sola algebra} era già stato
fatto in \texttt{\seqsplit{simmetrie\_correlatori\_trimero\_anello.tex}}: questo è
il primo che passa per il circuito quantistico reale su tutte e 81.

\textbf{Risultati.} I quattro zeri strutturali confermati esattamente (né un
quinto spurio, né uno dei quattro previsto che risulti in realtà non nullo).
Errori uniformemente piccoli su tutte le 81: Trotter medio $8.8\times10^{-4}$,
massimo $2.6\times10^{-3}$; statistico medio $1.4\times10^{-2}$, massimo
$3.7\times10^{-2}$. Risultato nuovo, non previsto dalla sola simmetria:
$C_{33}^{zz}(t{=}1.3)\approx0.9995$ — autocorrelazione quasi congelata (DM piccolo
rispetto a scambio/campo, $\sigma_3^z$ quasi conservato). I correlatori non banali
più intensi sono le coppie sito-1/sito-3 ($|C_{31}^{zx}|,|C_{13}^{xz}|\approx0.663$),
non le autocorrelazioni del sito 1 usate come esempio pedagogico principale.

\textbf{Come si usa.} Script standalone, $\sim4$--$5$ minuti (81 combinazioni
$\times$ 2 esecuzioni statevector $+$ 2 esecuzioni shots, cache di
transpilazione); salva \texttt{scan81\_results.npz}/\texttt{.json}, riusati da
\texttt{\seqsplit{generate\_figure\_scan81.py}} e dal terzo notebook della fase
(sotto) per evitare di ripetere il calcolo.

\textbf{Verifiche.} Nessun caso patologico isolato; rieseguito con successo anche
fuori da questo ambiente (macchina di Samuele, Windows/conda) dopo la correzione di
un percorso assoluto rimasto per errore nella prima versione consegnata.

\textbf{Conclusione.} Copertura completa del correlatore raggiunta: nessuna delle
77 combinazioni non-nulle nasconde un residuo sistematico o un errore mascherato.

\section*{Tre notebook della fase: \texttt{\seqsplit{correlazioni\_trimero\_anello\_simmetria.ipynb}}, \texttt{\seqsplit{correlazioni\_trimero\_anello\_esplorazione.ipynb}}, \texttt{\seqsplit{circuito\_correlazioni\_trimero\_anello\_tutte.ipynb}}}

Mirror per $N=3$ della struttura a tre notebook già usata per il dimero
(\texttt{\seqsplit{circuito\_correlazioni\_tutte.ipynb}}), eseguiti per intero (non
solo listati di codice). Tutta la logica pesante è riusata per import dai quattro
moduli sopra invece di essere riscritta.

\textbf{\texttt{\seqsplit{correlazioni\_trimero\_anello\_simmetria.ipynb}}} — mirror
computazionale di
\texttt{\seqsplit{simmetrie\_correlatori\_trimero\_anello.tex}}: verifica passo per
passo (non solo teoria) che $S_z^{tot}$ non è conservato, che $U_\text{naive}$
fallisce, che $U_\text{anello}$ commuta con $H$ (anche su 20 punti casuali), il
Lemma sui 9 operatori di sito, i quattro zeri strutturali a precisione macchina, gli
zeri a $t=0$ da time-reversal, e lo scan completo delle 81 combinazioni via formula
spettrale.

\textbf{\texttt{\seqsplit{correlazioni\_trimero\_anello\_esplorazione.ipynb}}} —
mirror della parte "circuito" del dimero: disegno del circuito, validazione su 6
casi rappresentativi, convergenza Trotter, un confronto shot-vs-Trotter. Nessuna
riderivazione delle simmetrie.

\textbf{\texttt{\seqsplit{circuito\_correlazioni\_trimero\_anello\_tutte.ipynb}}} —
mirror snello: misura diretta di tutte le 81 combinazioni via circuito, heatmap
$9\times9$ con i quattro zeri evidenziati, vista d'insieme nel tempo (griglia
$9\times9$, deliberatamente ridotta a $N=100$/8 punti di $t$, scelta dichiarata
esplicitamente), selettore manuale (niente \texttt{ipywidgets}, bug di rendering
già noto), analisi escursione/simmetria/spettro, confronto finale con lo scan a
shot finiti precomputato da \texttt{\seqsplit{scan81\_trimero\_anello.py}}.

\textbf{Nota metodologica.} Costruendo il terzo notebook è emerso un bug di
escaping in due punti distinti e indipendenti: uno script generatore che scrive il
testo delle celle in stringhe Python non-raw può consumare come sequenza di escape
(\texttt{\textbackslash b}, \texttt{\textbackslash a}) un \texttt{\textbackslash
beta}/\texttt{\textbackslash alpha} destinato a comparire letteralmente nel LaTeX
della cella — sia nello script generatore stesso (backslash da raddoppiare), sia
di nuovo nel codice generato se usa un f-string non raw (va reso \texttt{rf"..."}).
Errore silenzioso: non un traceback all'origine, ma un \texttt{ParseException} di
matplotlib mathtext a valle, con lettere mancanti nel titolo del grafico come unico
indizio.

\textbf{Come si usano.} Notebook, eseguiti per intero; richiedono i quattro moduli
sopra nella stessa cartella (e, per l'ultima sezione del terzo,
\texttt{scan81\_results.npz} già calcolato).

\textbf{Verifiche.} Tutte e tre le celle eseguite senza errori (verificato
iterando su tutti gli output di tutte le celle); risultati numerici coerenti con
quanto già misurato negli script standalone.

\textbf{Conclusione.} Chiudono, per il trimero anello, il mirror dei tre notebook
della fase correlazioni già completata per il dimero.

\section*{Conclusione generale}

La fase di correlazioni dinamiche per il trimero anello è arrivata a copertura
completa: tutte le 81 combinazioni misurate via circuito effettivo (non solo
classicamente), sia a statevector sia a shot finiti, con i quattro zeri
strutturali predetti dalla simmetria confermati esattamente e nessun bug residuo
di implementazione (a differenza del dimero, dove l'implementazione aveva rivelato
due bug reali) — merito dell'analisi di simmetria condotta \emph{prima} del
circuito, non dopo. Un risultato fisico non previsto dalla sola simmetria è emerso
dallo scan esaustivo: $C_{33}^{zz}(t{=}1.3)\approx0.9995$, autocorrelazione
del sito 3 quasi congelata al punto di lavoro adottato. Aperti per il seguito:
derivazione principiata (stile $R_1$ del dimero) del punto Trotter dimostrativo
$R_0$ del trimero; rumore di gate reale (Parte 2 della tesi); Fase 5 (VQE con DM)
per la catena aperta.

\end{document}